\honest{Variable Question Fallacies}{Eliezer Yudkowsky}{2008}

While writing the dialogue in which Albert and Barry argue about the falling tree, I sometimes lost the feel for why anyone would argue like that. To get it back, I would repeat: ``Either the falling tree makes a sound, or it does not!''

But ``P or not-P'' is not always safe with English sentences. ``This sentence is false'' is neither true nor false, and then there is ``Have you stopped beating your wife?'' A classical logician can keep the rule, for instance by saying the liar is not a sentence, but that takes subtlety.

The rule works only if P stands for exactly the same thing in both halves. If Albert::sound is not Barry::sound, the tree can make one and not the other. I borrow the ``::'' from C++, where it says which package's Sound you mean. \nb{Using a term in one sense at one point and in another sense at another is the fallacy of equivocation, which Aristotle listed among the fallacies that depend on language. The post does not use the old name.}

Many words ``are really functions of implicit variables supplied by context.'' There is no ``the left side'' of the street where your grocery store is, only your left as you travel in some direction. Programs that parse language find this hard. My second example: ``Martin told Bob the building was on his left.'' Whose left, Bob's or Martin's? \nb{That is the question of whom ``his'' refers to. The same ambiguity arises in ``near his house,'' with no ``left'' in the sentence.}

When two concepts share a name, or two events share one mental file, or one word is read in different contexts, reality seems to shift. That is how the algorithm feels from inside. \nb{The account comes from ``How An Algorithm Feels From Inside,'' which offers a guess about brain design and no evidence about brains.}

This confuses ``undergraduates (and postmodernist professors),'' who think an ambiguous sentence shows that reality is unstable. I name none of them. I supply one of their speeches myself: ``There is no fixed truth!'' Its refutation I leave ``as an exercise to the reader.''

I admit that my own draft said ``2 + 2 = X'' ``may have no fixed truth-value.'' In fact it has no truth value at all until X is fixed. This fallacy sneaks in even when you know better.
